\section{Categorical approach to the construction of distributive combinations}
\label{sec:construction}

% \todo[inline]{find better name for section}

While using \textsc{Frex} does not require knowledge about category theory, \textsc{Frex} uses many category-theoretic concepts internally. Therefore, the construction of distributive combinations is firstly category-theoretic. However we won't go into details here, as this would require a course on 2-categories and multi-categories.

\subsection{Challenge about constructing distributive combinations}
\label{sec:challenge}

Constructing generic free algebras and free extensions for arbitrary theories is not hard. In fact, \S\ref{sec:free_algebras} gives an explicit construction for the free algebra of any theory. However, the equivalence relation between terms in the quotient is not necessary effective/decidable. This is problematic since we've seen in \S\ref{sec:solving} that we need to be able to prove equalities in the free algebra to prove any equation with free algebras.


Free extensions also have a generic construction: \textcite{frex_origin,idris_frex} state that the free extension of an algebra $A$ by $X$ is the category-theoretic coproduct of $A$ with the free $\mathcal{A}$-algebra over $X$. However, there is no generic formula for the coproduct of two algebras for an arbitrary theory.

The difficulty of the construction is for it to be effective so that equivalence between two terms can be determined computationaly. This forbids the use of quotients, category theory pushouts etc...

\subsection{Initial result about the free algebra for the distributive combination}

% Let $\mathcal{T}$ be a theory. We say that $\la F_\mathcal{T}, \eta^\mathcal{T}, \bind^\mathcal{T} \ra$ is the free algebra construction for $\mathcal{T}$ if for every set of variables $X$, $\la F_\mathcal{T} X, \eta_X^\mathcal{T} : X \to F_\mathcal{T} X \ra$ is the free $\mathcal{T}$-algebra over $X$. $\bind^\mathcal{T}$ is given by the universal property of the free algebra: for every function $g : X \to \underline{A}$, $\bind^\mathcal{T} g : F_\mathcal{T} X \to A$ is a homomorphism.

% \medskip

A first result was given by Ohad Kammar (my internship supervisor) about the free algebra for the distributive combination of a theory over a commutative theory. A theory $\mathcal A$ is commutative iff for every $f : n \in \Sigma_\mathcal{A}$ and every $g : m \in \Sigma_\mathcal{A}$,
% \begin{equation}
%   f \bigl( g(x_{11}, \dots, x_{1m}), \dots, g(x_{n1}, \dots, x_{nm}) \bigr) = g \bigl( f (x_{11}, \dots, x_{n1}), \dots, f (x_{1m}, \dots, x_{nm}) \bigr)
%   \label{eq:commutativity_theory}
% \end{equation}
\begin{equation}
f \begin{pmatrix}
  g(x_{11}, \dots, x_{1m})\\
  \vdots\\
  g(x_{n1}, \dots, x_{nm})
\end{pmatrix} = g \begin{pmatrix}
  f \begin{pmatrix}
    x_{11}\\
    \vdots\\
    x_{n1}
  \end{pmatrix} & \cdots & f \begin{pmatrix}
    x_{1m}\\
    \vdots\\
    x_{nm}
  \end{pmatrix}
\end{pmatrix}
  \label{eq:commutativity_theory}
\end{equation}
For monoids, this states that \begin{gather*}
  (x \cdot y) \cdot (z \cdot w) = (x \cdot z) \cdot (y \cdot w) \qquad 0 \cdot 0 = 0 \qquad 0 = 0
\end{gather*} The theory of commutative monoids is therefore a commutative theory, fortunately. Note that every operator of arity $0$ and $1$ commutes with itself.

\begin{theorem}[Free algebra for the distributive combination, incomplete]
  Let $\mathcal A$ be a commutative theory, and $\mathcal M$ be any theory. Then there exists a distributive law of the free model monad induced by $\mathcal M$ over the free model monad induced by $\mathcal A$, and the free model monad induced by $\mathcal M \triangleright \mathcal A$ is the composite monad of the two free model monads.
  \label{thm:fail}
\end{theorem}

% \begin{theorem}[Free algebra for the distributive combination, incomplete]
% Let $\mathcal A$ be a commutative theory, and $\mathcal M$ be any theory. Suppose that we have $\lp F_\mathcal{A}, \eta^\mathcal{A}, \bind^\mathcal{A} \rp$ and $\lp F_\mathcal M, \eta^\mathcal{M}, \bind^\mathcal{M} \rp$ the free algebra constructions for $\mathcal A$ and $\mathcal M$.

% \noindent We can construct the free algebra construction $\lp F_\mathcal{\mathcal M \triangleright \mathcal A}, \eta^\mathcal{\mathcal M \triangleright \mathcal A}, \bind^\mathcal{\mathcal M \triangleright \mathcal A} \rp$ for the distributive combination $\mathcal M \triangleright \mathcal A$ using this data. The construction is effective, as defined in \S\ref{sec:challenge}. \todo[inline]{reformulate}
% \label{thm:fail}
% \end{theorem}

\noindent Intuitively, this means that given two free algebras for the additive and multiplicative part, composing them gives the free algebra for the distributive combination of the theories. This construction is effective in the sense of \S\ref{sec:challenge}.

Even though the construction itself was well-defined, the proof wasn't complete and that for a good reason: the statement is incorrect. The culprit is the absence of distributive laws under some hypothesis.

\subsection{Failure of the first result}

\textcite{nogo_distributive_law} gives multiple results about the absence of distributive laws under some hypothesis. We can regroup these results under 3 categories:
\begin{itemize}[$\triangleright$]
  \item theorems 3.5, 3.12 and 4.24 all forbid the existence of an idempotent operator in the multiplicative theory, on top of other conditions.
  \item theorem 4.7 forbids the existence of two distinct constants in the additive theory.
  \item theorem 4.17 imposes the abides equation $(a \ast b) \ast (c \ast d) = (a \ast c) \ast (b \ast d)$ to hold for every operator $\ast$ in the additive theory.
\end{itemize}

The second and third point are both taken care of by the hypothesis of commutativity of the additive theory in Theorem~\ref{thm:fail}. Two constants commute with each other if they are equal, and the abides equation is the commutativity hypothesis for $f = g = \ast$.

\medskip

However the first point is not accounted for. Under the current hypothesis, idempotent operators are allowed. Therefore, we need to strengthen our hypothesis.


\subsection{Corrected result}

We strengthen the hypothesis on $\mathcal{M}$ by requiring it to be an affine theory. A theory $\mathcal{M}$ is \emph{affine} if every axiom is an equation between terms where each variable occurs at most once. This allows weakening/deletion of variables, but not contraction. For example, $x, y \vdash x + y = y + x$ or $x \vdash x \cdot 0 = 0$ are affine equations, but $x \vdash x \vee x = x$ is not. Semigroups, monoids, groups and their commutative variants are all affine theories, whereas semilattices are not. This effectively covers the first point of the previous section.

\medskip

We now state the corrected theorem.

\begin{theorem}[Free algebra for the distributive combination]
Let $\mathcal A$ be a commutative theory, and $\mathcal M$ be an affine theory. Then there exists a distributive law of the free model monad induced by $\mathcal M$ over the free model monad induced by $\mathcal A$, and the free model monad induced by $\mathcal M \triangleright \mathcal A$ is the composite monad of the two free model monads.
\label{thm:correct}
\end{theorem}
 
The actual construction uses operads, translations from operads to theories, 2-categories and multicategories. Therefore, we won't present it nor the proof as is. We are gonna explain the gist of it.

First, start by noticing that a model for the distributive combination is made of a $\mathcal{M}$-model and a $\mathcal{A}$-model on the same carrier set (this ensures that the respective axioms hold) where moreover the operations of $\mathcal{M}$ distribute over the operations of $\mathcal{A}$ (this is a pullback in the category of categories). Let $X$ be a set of variables. The construction start by taking the free $\mathcal{M}$-algebra over $X$, $A_0 := F_\mathcal{M} X$. We then take the free $\mathcal{A}$-algebra over $\underline{F_\mathcal{M} X}$, $A_1 := F_\mathcal{A} \bigl( \underline{F_\mathcal{M} X} \bigr)$. This gives us our candidate $\mathcal{A}$-model. This step corresponds to the monadic composition, minus the multiplications. But we now need to define the semantics of the operators of $\mathcal{M}$ on the carrier $\underline{A_1}$. To do this, we "lift" the $\mathcal{M}$-operators so that the lifted operators on the new carrier distribute over the $\mathcal{A}$-operators. This is what gives the distributive law. Through some complex reasoning and with the help of the affine hypothesis, we prove that the $\mathcal{M}$-axioms hold. The commutativity hypothesis of $\mathcal{A}$ is used to prove the existence of an adjunction between multicategories, which is used to lift the $\mathcal{M}$-operators.

This construction is effective: it is solely built from applying the free algebra constructions and some 2-functors in well-chosen categories. The lifted operators can be defined concretely.

\subsection{An attempt at constructing the free extension for the distributive combination}

The natural next step is wanting to extend this result to free extensions. However the task is not so easy.

Not easy does not mean impossible. We managed to define a construction that uses pushouts and the free extensions of the component theories. While not effective because of the quotient nature of pushouts, this proves that it is indeed possible to use the free extensions of the composite theories to construct the one for the combined theory.

The construction of free extensions for distributive combinations is still ongoing work. We need to find a new pullback that is suited for extensions, and we need to find a new adjunction between the multi-category of extensions and the obvious forgetful functor.